\chapter*{Abstract}
\thispagestyle{plain}
\addcontentsline{toc}{chapter}{Abstract}

Multi-object tracking by detection is the standard way to obtain persistent object identities in video: a detector proposes scored boxes in every frame, and an online tracker associates them over time. Both components are governed by a small number of operating-point parameters, namely the detector's confidence threshold, suppression threshold, and input resolution, and the tracker's association and track-birth thresholds. These parameters are usually fixed once, for one detector and one type of scene, although scenes change within a single video and detectors are replaced, retrained, or quantized more often than trackers are retuned. This thesis studies whether training-free, causal control of these operating points can make tracking pipelines more accurate and more robust, and it measures where such control helps and where it does not.

The thesis first reviews real-time object detection and tracking, their paradigms, metrics, and datasets. It then presents two controllers that share the same principles: no learned weights, no labelled data at run time, one configuration frozen before any held-out evaluation, and paired bootstrap statistics for every comparison.

The first controller is scene-adaptive. It combines five inexpensive cues (detection density, tiny-object ratio, edge strength, low light, and blur), ranked against a label-free calibration set, into a scene complexity index that sets the detector's confidence threshold, suppression threshold, and resolution every ten frames in unmanned aerial vehicle video. On seventeen held-out VisDrone2019 sequences the validation-selected controller raised higher order tracking accuracy (HOTA) from 28.43 to 33.84 percent over a static default, at 38.98 processed frames per second on a Tesla T4 excluding image decoding, and it transferred to UAVDT without retuning (24.08 to 28.39 percent). A matched static operating point reached 33.93 percent, which shows that the gain comes from calibrating the operating point on aerial validation data rather than from frame-to-frame switching.

The second controller addresses what the first cannot observe: a change of the detector's confidence scale. It is a layer between a frozen detector and a frozen tracker that recalibrates itself from the detector's score stream with nested Otsu thresholds, decides whether the stream is clean or noisy, preserves the tracker's operating point on clean streams, and otherwise changes which candidates the tracker receives, their scores, and the tracker's thresholds through a declared host contract. One frozen configuration was evaluated on nine published trackers, four sources of detections, and three datasets. It raised HOTA by 0.61 points for two single-stage trackers, one of them predeclared and external; it restored trackers that halved detector scores had driven to zero HOTA to about 66; with a noisy detector it raised the multi-object tracking accuracy of two trackers by 6.0 and 7.9 points; and it changed HOTA by at most 0.06 points on five trackers with a low-score stage and well-matched detections. One external evaluation cell degraded significantly (1.52 HOTA points) because of false noisy decisions in a short sequence. The layer costs 2.95 ms per frame on a four-core processor.

Together, the two studies show that the operating point of a tracking pipeline should be calibrated before any adaptation is added, that adaptation should be judged against a matched static operating point, and that a score-driven, self-limiting layer can keep published trackers working when their detector changes, without claiming to improve every tracker.

\vspace{0.5cm}
\noindent\textbf{Keywords:} multi-object tracking, tracking by detection, detector operating point, confidence calibration, training-free control, unmanned aerial vehicles, Otsu thresholding.
